\item Identifica los elementos y los diferentes tipos de movimiento relacionados con la velocidad y aceleración y realiza experimentos sencillos.
\item Identifica y describe la presencia de fuerzas en interacciones cotidianas.
\item Experimenta e interpreta las interacciones de la fuerza y el movimiento, relacionados con las Leyes de Newton para explicar actividades cotidianas.
\item Experimenta e interpreta las interacciones de la fuerza y el movimiento relacionados con los principios de Pascal y de Arquímedes, para explicar actividades cotidianas.
\item Identifica algunos dispositivos de uso cotidiano en los cuales se aplica el Principio de Pascal (sistemas de frenos hidráulicos, elevadores y gatos hidráulicos) y de Arquímedes (flotación de barcos, submarinos y globos aerostáticos, entre otros); colabora en equipo para proponer actividades experimentales y resolver problemas sencillos relativos a las propiedades de los fluidos.
\item Analiza las características de la energía mecánica (cinética y potencial) y describe casos donde se conserva.
\item Relaciona al calor como una forma de energía y describe los motores que funcionan con energía calorífica, los efectos del calor disipado, los gases expelidos y valora sus efectos en la atmósfera.
\item Identifica saberes, prácticas y artefactos sobre el aprovechamiento de las diversas formas de energía renovables y no renovables, su empleo y origen en su comunidad (solar, eólica, hidráulica, geológica, mareomotriz, nuclear) y valora sus beneficios.
\item Realiza experimentos en donde se aproveche la energía del sol ya sea considerando las propiedades de la luz (energía solar) o las de la materia (convección).